\SourceLeaf{74}
\Needspace{8\baselineskip}
Man setze ferner

\[
\int f''(V - V^0)\,dp = W',
\]

\noindent die Integration so gemeint, dass bloss $p$ als variabel (d.\,i. $x, q, r$ etc. als constant) betrachtet werden, und

\[
dW' = V'dx.
\]

\noindent Dann ist, weil $W'$ kein $y$ enthält,

\[
fV' = 0, \qquad f'V' = 0, \qquad f''V' = f''(V - V^0),
\]

\noindent also auch

\[
{}'f(V - V^0 - V') = 0, \qquad f'(V - V^0 - V') = 0, \qquad f''(V - V^0 - V') = 0,
\]

\noindent und $V - V^0 - V'$ wird weder $y$ noch $p$ enthalten.

\Needspace{8\baselineskip}
Man setze ferner

\[
\int f'''(V - V^0 - V')\,dq = W'' \qquad \text{und} \qquad dW'' = V''dx,
\]

\noindent so wird ebenso

\[
\begin{aligned}f(V - V^0 - V' - V'') &={} f'(V - V^0 - V' - V'')\\ &={} f''(V - V^0 - V' - V'')\\ &={} f'''(V - V^0 - V' - V'') = 0,\end{aligned}
\]

\noindent und $V - V^0 - V' - V''$ enthält weder $y$ noch $p$ noch $q$, sondern bloss $x, r, s$ etc.

Ebenso enthält

\[
V - V^0 - V' - V'' - V''' \text{ bloss } x, s, t,
\]
\[
V - V^0 - V' - V'' - V''' - V^{\mathrm{IV}} \text{ bloss } x, t,
\]
\[
V - V^0 - V' - .\,.\,. - V^{\mathrm{V}} \text{ bloss } x.
\]

\noindent Es sei nun

\[
\int (V - V^0 - V' - .\,.\,. - V^{\mathrm{V}})\,dx = \Omega,
\]

\noindent so ist

\[
\int V dx = \Omega + W^0 + W' + W'' + .\,. + W^{\mathrm{V}},
\]

\noindent welches zu beweisen war.

